% problem-000165 7 镁铁氢化物晶体结构
\section*{7. 镁铁氢化物晶体结构}
2022年9⽉9⽇，国家航天局、国家原⼦能机构联合发布：我国科学家在嫦娥五号⽉壞样品中发现⼀种新矿物，命名为“嫦娥⽯”，英⽂名称Changesite-(Y)。嫦娥⽯是⼀种磷酸盐矿，属于陨磷钠镁钙⽯(Merrillite)族。

（图：）

\subsection*{7-1}
Merrillite 理想组成为 $\mathrm { C a _ { 9 } N a M g ( P O _ { 4 } ) } \div$ _{7}，晶体结构属三⽅晶系，晶胞参数 $\acute { \iota } a = 1 . 0 2 9 \mathrm { n m } , \ c =$ 3.687 nm。晶胞⽰意图⻅图7.1，其中(b)给出 $\mathrm { M g ^ { 2 + } }$ +在晶胞中的分布，可以看成沿�⽅向采⽤ABCABC⽅式排布。

\subsection*{7-1}
-1 写出晶胞中化学式的数⽬。

\subsection*{7-1}
-2 写出图7.1所⽰晶胞中所有 $\mathrm { M g ^ { 2 + } }$ +的分数坐标。（注意坐标轴的关系。要求按�值从⼩到⼤排列。）

\subsection*{7-1}
-3 计算 $\mathrm { M g ^ { 2 + } }$ 离⼦之间的最短距离。

\subsection*{7-2}
有趣的是，嫦娥⽯中含有稀⼟元素钇，这也是其英⽂名称后缀Y的由来。若Merrillite中所有$\mathrm { M g ^ { 2 } }$ +均被 $\mathrm { F e ^ { 2 + } }$ 取代，1/9 的 $\mathrm { { C a ^ { 2 + } } }$ 被 Y3+取代同时脱除 $\mathrm { N a ^ { + } }$ 以保持电荷平衡，即为嫦娥⽯，基本晶体结构和晶胞参数保持不变。

\subsection*{7-2}
-1 写出嫦娥⽯的化学式。

\subsection*{7-2}
-2 计算嫦娥⽯的密度（单位： $\mathrm { g c m ^ { - 3 } } )$ C

\subsection*{7-3}
初步研究表明，嫦娥⽯中含有 $^ { 3 } \mathrm { H e }$ ， $\mathrm { F e O } _ { 6 }$ 和 $\mathrm { P O } _ { 4 }$ 可能是其吸附位置。3He已⽤于中⼦探测，也很可能成为⼈类未来的能源。3He与中⼦作⽤产⽣质⼦和氚，写出核反应⽅程式。

（图：）
(a)

（图：）
(b)
图 7.1 陨磷钠镁钙⽯(Merrillite)的结构
